\subsection{The actual-exponent criterion before deck constantization}
\label{sec:nonconstant-deck-exponent}

\paragraph{Finding A45 (a geometrically cyclic pointed cover).}
Let $C/k$ be a smooth projective geometrically connected
curve of genus $g\ge2$ over a finite field of characteristic $p$,
and let $P\in C(k)$.
Let $\pi:D\to C$ be a geometrically connected finite etale cover
defined over $k$, such that $D_{\bar k}\to C_{\bar k}$ is cyclic
of odd degree $n>1$ prime to $p$. Assume that a point $P'\in D(k)$ lies over $P$.
The geometric deck group need not be constant over $k$. Put
\[
 B=\ker(\operatorname{Nm}:J_D\longrightarrow J_C)^0,
 \qquad M=\exp\bigl(B(k)[p^\infty]\bigr),
 \tag{ND.1}
\]
where the trivial group has exponent one. If
\[
                         n>(2M-1)^2,
 \tag{ND.2}
\]
then any distinct $Q\in C(k)$ whose point section collides with
that of $P$ forces a separable map
\[
 C_{\bar k}\longrightarrow E,\qquad
 \deg(C_{\bar k}/E)=p^a,\quad p\le p^a\le g-1,
 \tag{ND.2a}
\]
to an elliptic curve, totally ramified at both $P$ and $Q$.
Consequently (ND.2) separates $P$ from every distinct rational
point whenever $2\le g\le p$. In any genus it does so if no
map of the type (ND.2a) is totally ramified at $P$.
Only the rational point $P'$ is required; the entire fibre above
$P$ need not split over $k$. The exponent in (ND.1) is measured
over the original field $k$.

This extends A34 to nonconstant geometric deck groups and to
every characteristic, with the stated higher-genus obstruction. Whether it improves the numerical test for a specified
cover depends on the actual exponents before and after extending
constants. The primitive-conductor descent statement below records
an important case where those exponents are equal. No numerical
strictness is inferred just from the change of formulation.

\begin{proof}
\emph{1. Project the rational collision class before applying a deck map.}
Suppose that a distinct $Q\in C(k)$ has the same point section as
$P$. The pointed-cover implication supplies $Q'\in D(k)$ above
$Q$ with
\[
                       \delta=[Q'-P']\in J_D(k)[p^\infty].
 \tag{ND.3}
\]
This implication is the one used in A01 and A34; the geometric
monodromy of this cover is the cyclic group of order prime to $p$, so this is an
allowed prime-to-$p$ neighbourhood even though its arithmetic
Frobenius acts nontrivially on that group.

Choose a power $N$ of $p$ killing $\delta$ and an integer $r$
with $rn\equiv1\pmod N$. Define
\[
                b=\delta-r\pi^*\operatorname{Nm}(\delta).
 \tag{ND.4}
\]
All maps in this formula are defined over $k$, so $b$ is rational
over $k$ and is killed by $N$. Also $\operatorname{Nm}(b)=0$.
We verify that it lies in the connected norm kernel, without
assuming that the whole norm kernel is connected.

Over $\bar k$, choose a generator $\sigma$ of the geometric deck
group and write $S=1+\sigma+\cdots+\sigma^{n-1}$.
The norm--pullback identity is $\pi^*\operatorname{Nm}=S$.
One has $\operatorname{im}(\sigma-1)=B$: the image is connected
and lies in the norm kernel, whereas $n-S$ is divisible by
$\sigma-1$ in the integral group ring and restricts to $[n]$ on
$B$. Surjectivity of $[n]:B\to B$ proves the reverse inclusion.
Since $rn$ acts as the identity on the $N$-torsion class $S\delta$,
\[
             nb=(n-S)\delta\in B(\bar k),\qquad b=r(nb)\in B(\bar k).
 \tag{ND.5}
\]
Together with rationality, this proves $b\in B(k)[p^\infty]$.
The pullback term in (ND.4) is fixed by every geometric deck map.
Consequently
\[
              (\sigma-1)\delta=(\sigma-1)b,
              \qquad M(\sigma-1)\delta=0.
 \tag{ND.6}
\]
The class on the left need not itself be $k$-rational; its stated
annihilator follows from the rational class $b$, and requires no
extension of constants.

Let $m$ be its exact order, with $m=1$ for the zero class. Then
$m$ is a power of $p$ and $m\mid M$. On $D_{\bar k}$ there is a
function $f$ with
\[
 \operatorname{div}(f)
       =m(\sigma Q'+P'-Q'-\sigma P'),\qquad \deg f=d=2m.
 \tag{ND.7}
\]
The four points are distinct: the two base fibres are disjoint
and the deck action is free. If $m>1$ and $f$ were a $p$th power,
one $p$th of this divisor would be principal, contradicting the exact
order of $(\sigma-1)\delta$. If $m=1$, the simple poles already
exclude a $p$th power. Thus $f$ is separable over the perfect field
$\bar k$.

\emph{2. A small intermediate field.}
Set $K=\bar k(D)$ and
\[
 F=\bar k(\gamma f:\gamma\in\langle\sigma\rangle),
 \qquad e=[K:F],
\]
and let $Z$ be the smooth projective curve with function field
$F$. Poonen's Corollary~2.3, with Proposition~2.1 and Lemma~2.2,
applies because $F$ is generated by degree-$d$ functions; its
controlled-field integer is $d/e$. It gives
\[
                    e\mid d,\qquad g(Z)\le(d/e-1)^2.
 \tag{ND.8}
\]
The exact published locator is author p.~3
\cite{Poonen2007Gonality}. These extensions lie inside the
separable extension $K/\bar k(f)$. Since $D\to C$ is etale,
$g(D)=n(g-1)+1$. Inequality (ND.2) implies $n>(d-1)^2$, so $e=1$
is impossible. Hence $e>1$ divides $2p^b$ for some $b\ge0$, so
$\gcd(e,n)=1$. The resulting
map $v:D_{\bar k}\to Z$ is separable of degree $e$, and
$f=h\circ v$ for a function $h$ on $Z$.

\emph{3. The geometric deck group acts freely on $Z$.}
The field $F$ is stable under the group. Its action on $Z$ is
faithful, since a kernel would have order dividing both $n$ and
$e$. If a nontrivial stabilizer existed, it would contain a
subgroup of odd prime order $\ell$. This subgroup acts freely
on $D$ and preserves the fibre divisor over its fixed point
in $Z$. All its orbits in that divisor have length $\ell$ and
constant multiplicities, forcing $\ell\mid e$, a contradiction.
Thus the action is free. Riemann--Hurwitz gives
\[
                    g(Z)-1=n\bigl(g(Z/\langle\sigma\rangle)-1\bigr).
\]
The bound $0\le g(Z)<n$ forces both genera to be one. Therefore
\[
                   \deg\operatorname{Diff}(v)=2g(D)-2=2n(g-1).
 \tag{ND.9}
\]
This is the same geometric conjugate-field construction used in
A34 and recovered successor A.53 \cite[A.53]{ASrun}.

\emph{4. The two positive fibres force a small power of $p$.}
Suppose first that $P'$ and $\sigma Q'$ have the same image
under $v$. That fibre contains exactly these two positive
points: any other point would be an additional zero of $f$.
If the zero of $h$ has multiplicity $s$, both indices are
$m/s=e/2$. Equivariance gives the same index $e/2$ at $Q'$
and $\sigma P'$. Their pole images must coincide as well:
otherwise each would form a singleton degree-$e$ fibre and
have index $e$, a contradiction. Writing $x=v(P')$ and
$y=v(Q')$, these two fibre equalities give
$x=\sigma y$ and $y=\sigma x$. Thus $\sigma^2x=x$, contrary
to the free action, since $n>1$ is odd.

Therefore $P'$ and $\sigma Q'$ have distinct images in $Z$.
Each is the sole point in its zero fibre, and is totally
ramified with index $e$. Since its multiplicity in $f$ is $m$,
one has $e\mid m$. Thus $e$ is a nontrivial power of $p$.
The two deck orbits are disjoint because they lie above distinct
points of $C$, and each has $n$ elements. Every point in both
orbits has ramification index $e$ and different exponent at
least $e$. Indeed, for a local parameter expression
$u=a_et^e+\cdots$ the leading derivative vanishes because
$p\mid e$, so $\operatorname{ord}_t(du/dt)\ge e$; this derivative
is nonzero by separability. Together with (ND.9) this gives
\[
                   2ne\le2n(g-1),\qquad p\le e\le g-1.
 \tag{ND.10}
\]
For $g\le p$ this is already impossible, proving the stated
unconditional range.

\emph{5. Identify the exact elliptic obstruction in any genus.}
Let $E=Z/\langle\sigma\rangle$; it is a smooth genus-one
curve over $\bar k$, hence an elliptic curve after choosing
an origin. The fields satisfy
\[
 \bar k(E)=F^{\langle\sigma\rangle},\qquad
 \bar k(C)=K^{\langle\sigma\rangle},\qquad
 K=F\,K^{\langle\sigma\rangle}.
 \tag{ND.11}
\]
For the last equality, the subgroup of the Galois group of
$K/K^{\langle\sigma\rangle}$ fixing the compositum is the
kernel of its faithful action on $F$, and is therefore trivial.
Since $[F:\bar k(E)]=n$ and $[K:F]=e$, one obtains
$[\bar k(C):\bar k(E)]=e$. This extension is separable:
$K/F$ is separable and $F/\bar k(E)$ is etale. Thus it defines
a separable degree-$e$ map $C_{\bar k}\to E$.

The finite etale pullback $C_{\bar k}\times_E Z$ has function
field $K$ by (ND.11), so is connected and smooth and is precisely
$D_{\bar k}$. Ramification indices in $D_{\bar k}\to Z$ are
therefore those of $C_{\bar k}\to E$, because the other base
change map $Z\to E$ is etale. The total ramification at $P'$
and $\sigma Q'$ descends to total ramification at $P$ and $Q$.
Together with (ND.10), this proves the obstruction (ND.2a)
and all asserted sufficient cases.
\end{proof}

\paragraph{Integral semilinear descent for a primitive conductor.}
Let $p$ be a prime, let $n>1$ be prime to $p$, and let
$a\in(\mathbf Z/n)^\times$ have order
$f$.  Use the full cyclotomic algebra
\[
 \mathcal O_n=\mathbf Z_p[Z]/\Phi_n(Z),\qquad
 \tau(Z)=Z^a,\qquad R_n=\mathcal O_n^{\langle\tau\rangle}.
 \tag{SD.1}
\]
This algebra can be a product of unramified discrete valuation
rings.  It is not to be replaced by one chosen factor unless
that factor is preserved by $\tau$.
Let $M$ be a finite $\mathcal O_n$-module killed by a power of
$p$, with an additive automorphism $F$ satisfying
\[
 F(cm)=\tau(c)F(m),\qquad F^f=1.
 \tag{SD.2}
\]
Then the natural map and its exponent consequence are
\[
 \mathcal O_n\otimes_{R_n}M^F\ \xrightarrow{\ \sim\ } M,
 \quad c\otimes m\longmapsto cm,
 \qquad \exp(M)=\exp(M^F).
 \tag{SD.3}
\]
The exponent of the zero group is one.  No freeness of $M$,
and no assumption that $p$ is prime to $f$, is required.

\begin{proof}
Since $p\nmid n$, $Z^n-1$ has distinct roots modulo $p$.
Choose a finite unramified extension $S/\mathbf Z_p$ in which
these roots all split; one obtains it by lifting the distinct
roots over $\mathbf F_{p^h}$ with $n\mid p^h-1$.
The resulting faithfully flat base change identifies
\[
 \mathcal O_n\otimes_{\mathbf Z_p}S
                    \simeq\prod_{j\in(\mathbf Z/n)^\times} S.
 \tag{SD.4}
\]
The automorphism $\tau$ acts by multiplication by $a$ on the
primitive-root indices.  Every orbit has exactly $f$ elements:
$a^ij=j$ with $j$ a unit implies $a^i=1$ modulo $n$.
Taking invariants commutes with this flat base change, since
invariants are the kernel of $\tau-1$ on the finite free
$\mathbf Z_p$-module $\mathcal O_n$.  Hence
\[
 R_n\otimes_{\mathbf Z_p}S
       \simeq\prod_{(\mathbf Z/n)^\times/\langle a\rangle}S.
 \tag{SD.5}
\]
It follows, by faithfully flat descent, that
$\mathcal O_n/R_n$ is finite etale and faithfully flat of
rank $f$.  More explicitly, the map
\[
 \mathcal O_n\otimes_{R_n}\mathcal O_n
       \longrightarrow\prod_{i=0}^{f-1}\mathcal O_n,
 \qquad c\otimes d\longmapsto(c\tau^i(d))_i
 \tag{SD.6}
\]
is an isomorphism: after the faithfully flat base change $S$
it is the bijection of coordinates supplied by a free cyclic
permutation on each orbit.  This is the finite etale
$\mathbf Z/f$-torsor property, including when $p$ divides $f$.

The powers of $F$ in (SD.2) therefore give an actual descent
datum for the $\mathcal O_n$-module $M$ along $R_n\to\mathcal O_n$.
The cocycle condition is exactly $F^f=1$ and the semilinear
composition law.  Faithfully flat descent for modules, in its
unrestricted module form, identifies the descended module with
$M^F$ and gives (SD.3); the exact primary statement is the
Stacks Project, Proposition~35.3.9, Tag~\texttt{023N}
\cite{Stacks023N}.  In particular this step applies to finite
torsion modules and does not average by $1/f$.
If $p^e$ kills $M^F$, it kills the displayed tensor product and
therefore $M$.  Conversely $M^F$ is a subgroup of $M$.
This proves the equality of exponents.
\end{proof}

\paragraph{Application to geometric cyclic covers.}
Let $\pi:D\to C$ be a geometrically connected etale cover of
smooth proper curves over a finite field $k$ of characteristic
$p$, geometrically Galois with cyclic deck group of odd order
$n$ prime to $p$.  Choose a geometric generator $\sigma$ and let $F$ be
arithmetic Frobenius on geometric points.  Write
\[
                 F\sigma F^{-1}=\sigma^a,
       \qquad f=\operatorname{ord}_n(a),\qquad k_s=\mathbf F_{q^s}.
 \tag{SD.7}
\]
Thus $k_f$ is the minimal constant extension defining every
deck transformation.  Put
$B=\ker(\operatorname{Nm}_\pi:J_D\to J_C)^0$.
Because $n$ is a unit in $\mathbf Z_p$, the group algebra
decomposes into its cyclotomic factors:
\[
 \mathbf Z_p[\langle\sigma\rangle]
       \simeq\prod_{d\mid n}\mathcal O_d.
 \tag{SD.8}
\]
Indeed the reductions of the cyclotomic factors of $Z^n-1$
are pairwise coprime.  If $e_d$ is the resulting idempotent,
write
\[
 \mathscr T_d=e_d\bigl(B(\overline k)[p^\infty]\bigr),
              \qquad d\mid n,\quad d>1.
 \tag{SD.9}
\]
These are the primitive conductor-$d$ torsion modules.  They
are stable under $F$, since raising a primitive root to $a$
preserves its conductor.  The trivial conductor is absent
from $B[p^\infty]$, and
$B(\overline k)[p^\infty]=\bigoplus_{d\mid n,\ d>1}\mathscr T_d$.
One can see this directly from
$\pi^*\operatorname{Nm}_\pi=\sum_{i=0}^{n-1}\sigma^i$:
the average of this sum projects onto conductor one, and
$[n]$ is invertible on all $p$-primary groups.

These modules can also be realized by actual subvarieties,
without replacing a factor by an arbitrary isogenous variety.
Fourier orthogonality on the $n$th roots of unity gives
\[
 \begin{aligned}
 p_d&=ne_d=\sum_{j=0}^{n-1}c_d(j)\sigma^j,
                       &p_d^2&=np_d,\\
 c_d(j)&=\sum_{\substack{\zeta^d=1\\\operatorname{ord}(\zeta)=d}}
                         \zeta^{-j}\in\mathbf Z.
 \end{aligned}
 \tag{SD.9a}
\]
The coefficients are integers because they are rational
algebraic integers.  The projector is fixed by Frobenius,
which permutes the primitive characters.  Thus
$B_d=\operatorname{im}(p_d:B\to B)$ is defined over $k$.
On $B_d$ the map $p_d$ is $[n]$; since $p$ is prime to $n$ this shows
$B_d(\overline k)[p^\infty]=\mathscr T_d$.
Indeed both inclusions follow by applying $p_d$ and the
inverse of $[n]$ on the finite cyclic group generated by a
$p$-primary point.  Hence the following exponent statements
refer to these canonical actual conductor subvarieties.

For each $d>1$ put $f_d=\operatorname{ord}_d(a)$ and
\[
                      M_d=\mathscr T_d^{F^{f_d}}.
\]
This is finite, since it is a subgroup of $B(k_{f_d})$.
On $\mathcal O_d$, $F$ is semilinear for the automorphism
$Z\mapsto Z^a$ of order $f_d$, and $F^{f_d}=1$ on $M_d$.
The preceding lemma gives the exact equality
\[
       \exp\bigl(\mathscr T_d^{F^{f_d}}\bigr)
                    =\exp\bigl(\mathscr T_d^F\bigr).
 \tag{SD.10}
\]
Thus a primitive conductor factor has the same rational
$p$-primary exponent over $k$ as over its own minimal deck
constantization field.  In particular, when $n$ is prime,
the whole Prym has this equality between $k$ and $k_f$.
It also holds for the primitive conductor-$n$ factor of any
odd degree, including a prime power.

For the whole Prym of a composite-degree cover one must retain
the distinction between $f_d$ and $f$.  If $f_d<f$, the module
$\mathscr T_d^{F^f}$ is acted on by $F$ with $F^f=1$, but
$F^{f_d}$ need not be the identity.  The coefficient action
has already become trivial at $f_d$; its remaining action on
the module is an ordinary linear action.  Consequently
(SD.2) with exponent $f_d$ is not a hypothesis on this larger
module.  What descent identifies is exactly
\[
 \bigl(\mathscr T_d^{F^f}\bigr)^{F^{f_d}}
        =\mathscr T_d^{F^{f_d}},
 \tag{SD.11}
\]
whose exponent is given by (SD.10).  It gives no equality
between the exponent of $\mathscr T_d^{F^f}$ and that of
$\mathscr T_d^F$.  The full exponent is the maximum of the
conductor exponents, so proper-conductor growth is precisely
the part not excluded by this argument.


\paragraph{An actual improvement before the deck field: degree thirty-nine.}
Let $k=\mathbf F_4=\mathbf F_2(u)$, with $u^2+u+1=0$, and put
\[
 C/k:\quad Y^2+X^2Y=uX^5+uX^4+X^3+uX,
                       \qquad W=\infty.
 \tag{ND.12}
\]
There is a geometrically connected etale cyclic cover
$D\to C$ of degree $39$, with a $k$-point over $W$, whose
minimal deck field is $K=\mathbf F_{256}$. For its whole Prym
$B$ the exact exponents are
\[
       \exp B(k)[2^\infty]=2,
             \qquad \exp B(K)[2^\infty]=16.
 \tag{ND.13}
\]
Thus A45 applies over $k$, since $39>9$, whereas the actual
constant-deck test of A34 for this same cover over $K$ fails,
since $39\le961$. The whole new factor is ordinary of dimension
$38$, and its primitive conductor-$39$ factor already has
nonzero rational two-primary torsion over $k$. This is a
strict improvement for a specified pointed cover and field;
it asserts no absence of other successful covers.

\begin{proof}
\emph{1. The twist and the exact stable quotient.}
The curve (ND.12) is the quadratic twist of A27's model
$C_0:Y^2+X^2Y=uX^5+X^3+uX$. Adding $uX^4$ adds the constant
$u$ to the Artin--Schreier equation, and
$\operatorname{Tr}_{\mathbf F_4/\mathbf F_2}(u)=1$.
The genus is two and the geometric two-rank is one, as for
$C_0$. The point $W$ is rational and $K_C\sim2W$.

Use the field and integer encoding of A32:
\[
 K=\mathbf F_2[z]/(z^8+z^4+z^3+z+1),
       \quad u=188,\quad b=92,
       \quad b^2+b=u,\quad b^4+b=1.
 \tag{ND.14}
\]
The isomorphism to $C_0$ over $K$ is
$(X,Y)\mapsto(X,Y+bX^2)$. It fixes $W$.
The checked A27 polynomial gives $\#J_C(k)=18$, and A32 gives
$\#J_C(K)=65520=2^4\cdot3^2\cdot5\cdot7\cdot13$.
Write $G=J_C(K)$, identified with $J_{C_0}(K)$ by this map.
The exact reduced Mumford representatives
\[
 \Delta_9=(X^2+189,\,189X+80),\qquad
 \Delta_{13}=(X^2+80X+237,\,106X+207)
 \tag{ND.15}
\]
on $C_0$ have orders nine and thirteen. On the twist,
$\Delta_9=(X^2+u+1,(u+1)X+u)$ is already $k$-rational.
The finite computation verifies that these generate the entire
respective primary subgroups of $G$ and that the original
four-power Frobenius acts on them by $8$ modulo nine and $5$
modulo thirteen. The twist changes that Frobenius to its
negative on the Jacobian; the two actions are therefore $1$
modulo nine and $8$ modulo thirteen.

Define $a_9$ and $a_{13}$ by
\[
 [7280]Q=a_9(Q)\Delta_9,\qquad
 [5040]Q=a_{13}(Q)\Delta_{13},
\]
and let $a:G\to\mathbf Z/39$ be given by
$a\bmod3=a_9\bmod3$ and $a\bmod13=a_{13}$.
These are homomorphisms, and $a$ is surjective: multiplication
by the displayed cofactors kills the other primary factors
and is invertible on the selected primary subgroup.
For the twist Frobenius $\Phi$,
\[
                a(\Phi Q)=34a(Q)\pmod{39},
                        \operatorname{ord}_{39}(34)=4.
 \tag{ND.16}
\]
The action on the conductor-three quotient is trivial.
Its actions modulo thirteen and modulo thirty-nine both have
order four.

The finite etale isogeny $\Phi^4-1:J_C\to J_C$ is defined
over $k$ and has geometric kernel $G$.
Its subgroup $H=\ker a$ is Frobenius stable by (ND.16).
The quotient isogeny $J_C/H\to J_C$ is therefore defined
over $k$, has cyclic geometric kernel $G/H$ of order $39$,
and has deck Frobenius $34$. Pull it back through
$C\to J_C$, $R\mapsto[R-W]$. The result is geometrically
connected by Milne's full Abel--Jacobi fundamental-group
statement, Proposition~9.1 \cite{MilneJV}; the isogeny cover
is connected and etale. The zero point gives a $k$-point
over $W$. This constructs $D$ and proves that its minimal
deck field is precisely $K$.

\emph{2. Complete finite character data.}
Over $K$ this is the constant Lang quotient of $G$ by $H$.
For $c_j=\#\{R\in C(K):a([R-W])=j\}$, exact enumeration gives
\[
\begin{aligned}
 (c_0,\ldots,c_{12})&=(14,7,6,6,6,7,4,6,7,4,6,6,5),\\
 (c_{13},\ldots,c_{25})&=(15,7,6,6,6,5,6,6,5,6,6,6,7),\\
 (c_{26},\ldots,c_{38})&=(15,5,6,6,4,7,6,4,7,6,6,6,7).
\end{aligned}
 \tag{ND.17}
\]
The sum is $256$. Reducing the labels modulo three gives
$(74,91,91)$; reducing modulo thirteen gives exactly the A32
table (OR.9). The four rational points of (ND.12) over $k$ are
\[
              W,\quad(0,0),\quad(u,u+1),\quad(u,0).
\]
Their conductor-three labels are respectively $0,0,1,2$.
These are also their labels for the conductor-three cover
over $k$: extending constants from $k$ to $K$ multiplies a
degree-one Frobenius label by four, which is one modulo three.

For any nontrivial character $\chi$ of the indicated cyclic
quotient, the genus-two character polynomial at $W$ is
\[
         L_\chi(T)=1+c_\chi T+qT^2,
             \qquad c_\chi=\sum_R\chi(a([R-W])).
 \tag{ND.18}
\]
Here $q$ is the size of the field over which the deck group is
constant; this is the exact formula already proved in the
recovered A1 B.2 and used in A29 and A32. For conductor three
it gives $c_\chi=1$ over $k$, and $c_\chi=-17$ over $K$.
Thus $L_\chi(1)$ is respectively $6$ and $240$.
Since two is primitive modulo three, the integral module
argument (BE.8) applies to this genus-two degree-three cover.
It gives
\[
              \exp B_3(k)[2^\infty]=2,
                 \qquad \exp B_3(K)[2^\infty]=16.
 \tag{ND.19}
\]
Here the degree-thirteen quotient $D\to D_3$ identifies the
degree-three Prym's two-primary module with the actual $B_3$:
its pullback and norm have composite $[13]$ in both directions
on this conductor, and $[13]$ is invertible on every two-primary
group.
For conductor thirteen, the A32 character values give odd
$\#B_{13}(K)$, hence no rational two-primary torsion over
either field.

\emph{3. Exact exponent of the conductor-thirty-nine factor.}
Let $Z$ denote a primitive thirty-ninth root formally, and
reduce $257+\sum c_jZ^j$ modulo $\Phi_{39}(Z)$. The result is
\[
\begin{aligned}
 L_{39}(1)=2\bigl(&128+Z+Z^4-Z^6+Z^7-Z^9-Z^{12}\\
                &+Z^{14}+Z^{17}-Z^{18}+Z^{20}-Z^{21}\bigr).
\end{aligned}
 \tag{ND.20}
\]
Write $\beta$ for the reduction modulo two of the expression
in parentheses. It is a unit modulo $\overline\Phi_{39}$.
For a direct coefficient certificate, the polynomials
\[
\begin{aligned}
 U&=Z^{20}+Z^{17}+Z^9+Z^7+Z^6+Z^4,\\
 V&=Z^{17}+Z^{14}+Z^{11}+Z^{10}+Z^9+Z^8
                         +Z^6+Z^5+Z^2+Z+1
\end{aligned}
\]
satisfy $U\beta+V\overline\Phi_{39}=1$ in $\mathbf F_2[Z]$.
Thus in the full finite etale algebra
$\mathcal O_{39}=\mathbf Z_2[Z]/\Phi_{39}$, the value (ND.20)
is two times a unit. Its norm has two-adic valuation
$\varphi(39)=24$. Character factorization and the norm--pullback
isogeny therefore give
\[
                       v_2\bigl(\#B_{39}(K)\bigr)=24.
 \tag{ND.21}
\]

Moreover (ND.20) says $c_\chi\equiv1\pmod2$ for every
primitive character. Each corresponding coherent $H^1$
summand has dimension one, because the base genus is two.
As in A32, the reduced character polynomial is the determinant
of linear $256$-power Frobenius on that summand. It is $1+T$;
thus the scalar is one on every summand. The functorial
Artin--Schreier stable-part comparison makes $B_{39}$ ordinary
of dimension $24$ and all its geometric etale two-torsion
rational over $K$. The exact cited comparison inputs are
Deligne, Theorem~2.2(b) and formula~(2.4.1), with the hypotheses
of Section~2.1 \cite{DeligneModP}, and Katz,
Propositions~2.2.4--2.2.5 \cite{KatzCongruence}, as specified
in A32. The deck group is constant over $K$, so its distinct
odd-order characters split the coherent space; there is no
uncontrolled Jordan block between these one-dimensional lines.

Consequently $B_{39}(K)[2]$ already has order $2^{24}$.
By (ND.21) this is the entire two-primary group, proving
\[
                        \exp B_{39}(K)[2^\infty]=2.
 \tag{ND.22}
\]
Primitive semilinear descent, (SD.10), gives the same exponent
over $k$, since the conductor-thirty-nine action in (ND.16)
has order four. The conductor idempotents have odd denominators
and split the two-primary group of the actual whole Prym into
its conductor-$3$, $13$, and $39$ components. Taking the maximum
of (ND.19), (ND.22), and the conductor-thirteen exponent one
proves (ND.13). Each of the three geometric components is
ordinary, of dimensions $2$, $12$, and $24$, respectively;
hence $B$ is ordinary of dimension $38$. The cover itself has
genus $40$ and two-rank $39$.

Finally $g=2=p$ and $39>(2\cdot2-1)^2$, so A45 separates
the section of $W$ over the original field. The comparison
over the full deck field uses the actual larger exponent
$16$ and fails the inequality. This proves the stated strict
comparison for this cover.
\end{proof}

\paragraph{Finite computation scope.}
The table calculation \texttt{probe\_A27\_c39.py} uses the
retained, independently checked A32 point list and its binary
linear-interpolation Jacobian arithmetic. It recomputes the
three-primary projections, their cyclic order-nine structure,
the Frobenius actions, the twist's four rational points, the
new thirty-nine labels, the cyclotomic reductions, and the
displayed Bezout identity. Its JSON retains all $256$ point
images and the exact data above. This use of retained arithmetic
is not described as a newly independent Jacobian implementation.


The separate checker \texttt{verify\_A45\_conductor39\_cantor.py}
uses the previously checked extended-GCD Cantor arithmetic,
imports none of the linear-interpolation producer, and freshly
enumerates every point and both primary projections. It verifies
all $256$ new quotient labels, both complete cyclic subgroups,
the twisted Frobenius permutation of every point, the four
$k$-points, and the two exact cyclotomic unit identities.
Both implementations pass and agree on every point label.
A further rerun of that checker reproduces its JSON exactly;
that rerun is not counted as a third independent implementation.
